% ch-ref-numerics.tex -- reference: the Numerics and Restart blocks
\chapter{Numerics and restart}
\label{ch:ref-numerics}

\key{Numerics} holds every setting that decides \emph{how} the equations are
solved rather than what they are: the time step, the Newton iteration of the
height equation, and the two linear solves. \key{Restart} says where a
continued run picks up from. This is a reference chapter --- it gives the keys,
their ranges and their immediate meaning, and points at the numerical-method
part for the reasoning, which is where you should go when a run will not
converge or will not move.

\section{\key{Numerics}}
\label{sec:ref-numerics:numerics}

Required, exactly once. These five keys govern the outermost loop: how long the
run lasts and when it decides that nothing is happening any more. The
steady-state test measures the relative change of thickness \emph{per unit
time}, multiplied by the total time span --- ``how much the shape would still
move if the run continued at the present rate'' --- rather than the change per
step, so that a collapsed time step cannot be mistaken for a steady state. See
\cref{ch:timeintegration}.

\begin{keytable}{Keys of \key{Numerics}\label{tab:numerics-keys}}\key{InitialDt} & double & \reqd & Initial time step in seconds, positive. Clamped down to the smallest output interval if it is larger. \\
\key{EndTime} & double & \val{0} & Physical end time in seconds; the loop runs while $\mathrm{StartTime}\le t\le\mathrm{EndTime}$. Positive in practice --- see the warning below. With a restart it is still an absolute time, so it must exceed \key{Restart.RestartTime}. \\
\key{MaxSteps} & int & \val{0} & Stop after this many steps, whatever \key{EndTime} says. At least 0; \val{0} sets no limit. Counted in 64 bits, up to $2^{63}-1$; a restart counts on from the steps already taken. \\
\key{SteadyStateCriterion} & double & \val{1e-3} & Threshold on (relative thickness change per unit time) $\times$ ($\mathrm{EndTime}-\mathrm{StartTime}$), non-dimensional. Any value; $\le 0$ disables the test. The span falls back to 1 when $\mathrm{EndTime}\le\mathrm{StartTime}$. \\
\key{SteadyStateSteps} & int & \val{10} & How many consecutive steps the criterion must hold, range $[1:1000000]$. The test is \code{count > SteadyStateSteps}, so the run needs one step more than the number given. \\
\end{keytable}

\begin{warn}[\key{EndTime} defaults to 0, which is not ``run for ever'']
  With \key{EndTime} left at its default the run stops after a single step, and
  the steady-state test falls back to a time span of 1. There is no setting
  that means ``unlimited''. \textbf{Always set \key{EndTime} explicitly.} The
  other two ways a run can stop are the dry-out test
  (\key{Physics.Evaporation.EndConcValue}) and the steady-state test above.
\end{warn}

\section{\key{Numerics.TimeStepControl}}
\label{sec:ref-numerics:tsc}

Required block. The time step is adaptive only when \key{Enable} is true; the
step is decided by the height equation, and the concentration equation follows
it. The error estimate comes from comparing two consecutive time differences,
and the step shrinks when the estimate exceeds \key{Criterion}. See
\cref{ch:timeintegration}.

\begin{keytable}{Keys of \key{Numerics.TimeStepControl}\label{tab:tsc-keys}}\key{Enable} & bool & \val{false} & Make the time step adaptive. When false, $\Delta t$ stays at \key{InitialDt}. \\
\key{MaxDtFactor} & double & \val{10000} & Upper bound on $\Delta t$ as $\mathrm{MaxDtFactor}\cdot\Delta x_{\min}^4$, the fourth-order-diffusion scaling, in non-dimensional time. Positive. Applied whether or not \key{Enable} is true. \\
\key{MinDt} & double & \val{1e-15} & Floor on $\Delta t$ in non-dimensional time, positive. The run aborts when a smaller step is requested. Applied whether or not \key{Enable} is true. \\
\key{Criterion} & double & \val{0.01} & Tolerated local truncation error per step, non-dimensional, positive. \\
\key{ScaleFloor} & double & \val{0.0} & Floor on the weight that normalizes the error, as a fraction of the global maximum thickness, range $[0:1]$. \val{0} divides by the local thickness, so thin precursor regions dominate; \val{1} normalizes by the global maximum only. \\
\key{KeepSteps} & int & \val{50} & How many consecutive successful steps are required before $\Delta t$ may grow again, at least 1. \\
\key{ChangeFactor} & double & \val{0.5} & Factor $\Delta t$ is multiplied by when it shrinks, and divided by when it grows. Range $(0:1)$. \\
\end{keytable}

\begin{note}[What the \key{Enable} switch does and does not cover]
  With \key{Enable = false}, \key{Criterion}, \key{ScaleFloor},
  \key{KeepSteps} and \key{ChangeFactor} are read but never applied ---
  whereas \key{MaxDtFactor} and \key{MinDt} still are. A fixed step is
  therefore still capped from above and floored from below. The \val{"CZ"}
  model applies neither of those two: its step is capped by the output
  intervals alone, and floored at the built-in \val{1e-15}. A step that
  does not converge --- the well-mixed Newton iteration, which includes a
  correction that stagnates far from convergence
  (\cref{sec:convergence:stagnation}), or a linear solve
  of the film or of the concentration in either model --- cannot be retried
  at a shorter one, so it stops the run with that reason. With
  \key{Enable = true} such a step is cut and solved again.
\end{note}

\section{\key{Numerics.HeightEquation}}
\label{sec:ref-numerics:height}

Required block. The height equation is non-linear, because its mobility
depends on the unknown, so each step is solved by an outer Newton--Kantorovich
iteration. These five keys control that iteration. They are the
settings most worth understanding before you change anything else: the single
most common cause of a run that crawls is a convergence test that cannot be
satisfied, which collapses the time step. \Cref{ch:convergence} explains the
mechanism in full, and \cref{sec:convergence:floor} explains why the
correction cannot fall below a floor, and when that floor is set by the mesh
and the step size.

\begin{keytable}{Keys of \key{Numerics.HeightEquation}\label{tab:heighteq-keys}}\key{NewtonEpsilon} & double & \val{1e-10} & Absolute convergence threshold on the Newton--Kantorovich correction, non-dimensional, positive. Setting it below the floor of the correction makes every step ``fail'' and collapses the time step. \\
\key{NewtonEpsilonRel} & double & \val{1e-6} & Relative threshold: how far the correction must fall below the step's \emph{first} correction. Range $[0:1]$; $\le 0$ disables it, leaving the absolute test alone. Independent of solution scale and of $\Delta t$. \\
\key{NewtonMaxIterations} & int & \val{30} & Iteration cap for the Newton solve, at least 1 in practice. Exceeding it counts as a failed step, which shrinks $\Delta t$ when the adaptive step is on and stops the run at a fixed step; so does a correction that stagnates far from convergence (\cref{sec:convergence:stagnation}). \\
\key{NewtonMinIterations} & int & \val{1} & Minimum number of Newton iterations per step, range $[1:100]$. \val{1} lets the convergence test alone decide; \val{6} makes every step take at least six. \\
\key{JacobianUpdateInterval} & int & \val{1} & Modified-Newton interval: how often the Jacobian is rebuilt, range $[1:100]$. \val{1} rebuilds it every iteration. \\
\end{keytable}

\begin{tip}[Prefer the relative test to a smaller \key{NewtonEpsilon}]
  When a run stalls, lowering \key{NewtonEpsilon} is the wrong direction: where
  the correction has a floor proportional to $\Delta t$, a smaller absolute
  threshold feeds back into a smaller step, which lowers the floor, and the
  run locks onto the step where the two just meet. Leave the absolute test
  where it is and let \key{NewtonEpsilonRel} do the work
  (\cref{sec:convergence:floor}).
\end{tip}

\begin{warn}[The \val{"CZ"} model ignores all five of these keys]
  With \key{Physics.ConcField.Model = "CZ"} the height equation is solved
  semi-implicitly and never calls the corresponding setters, so
  \key{NewtonMaxIterations}, \key{NewtonMinIterations}, \key{NewtonEpsilon}, \key{NewtonEpsilonRel} and
  \key{JacobianUpdateInterval} have no effect and the built-in defaults of the
  Newton solver are used instead. The same model also ignores
  \key{Physics.DisjoiningPressure} and the whole \key{Physics.Source} block.
\end{warn}

\section{\key{Numerics.ConcEquation}}
\label{sec:ref-numerics:conceq}

Required block with \textbf{no keys of its own}. Its sole purpose is to hold
the \key{LinearSolver} block of the concentration equation. It must be present
even when no concentration equation is solved --- with \key{Model = "None"},
for instance, the block is still required, and its contents are still ignored.

\section{\key{Numerics.Convection}}
\label{sec:ref-numerics:convection}

Optional, at most once. It chooses the face interpolation of the convective
terms (\cref{sec:disc:convection}).

\begin{yzntable}{Keys of \key{Numerics.Convection}\label{tab:ref-numerics:convection}}{@{}L{0.20\textwidth} L{0.08\textwidth} L{0.14\textwidth} L{0.50\textwidth}@{}}
\textsf{\bfseries Key} & \textsf{\bfseries Type} & \textsf{\bfseries Default} & \textsf{\bfseries Meaning}\\ \midrule
\key{Concentration} & string & \val{"UTOPIA"} & The advection of $c$: in $x$ and $y$, and in \val{"CZ"} also in $\zeta$. \\
\key{Height} & string & \val{"Koren"} & The upwinded fluxes of the well-mixed film-height equation: the in-plane gravity, the centrifugal term and the Marangoni flux. Refused with \val{"CZ"}, whose film-height equation upwinds nothing. \\
\end{yzntable}

Both take one of \val{"UTOPIA"} (third-order upwind), \val{"KK"}
(Kawamura--Kuwahara: the same with three times the dissipation),
\val{"Koren"} and \val{"VanLeer"} (limited, bounded), \val{"Central4"},
\val{"QUICK"}, \val{"Upwind1"}, \val{"Upwind2"} and \val{"Central2"}; any
other name is refused with the list. The schemes in use are written to the
log and to \file{JobInfo.txt}. \Cref{sec:verification:schemes} gives the comparison
the defaults were chosen by: \val{"Koren"} for the film height because it
keeps the face height between its neighbors and survives steep fronts over
a thin precursor, \val{"UTOPIA"} for the concentration because it is the
most accurate on smooth profiles. For a concentration with sharp fronts,
\val{"Koren"} keeps $c$ within its bounds at the cost of accuracy at smooth
extrema. \val{"Central2"} and \val{"Central4"} dissipate nothing: where the
model itself is unstable at the shortest resolved waves --- in \val{"CZ"}, a
film stratified with $\sigma_c\,\partial c/\partial\zeta > 0$ and a weak
diffusivity --- they let those waves grow
(\cref{sec:verification:schemes}).

\paragraph{The Courant check of the limited schemes.} \val{"Koren"} and
\val{"VanLeer"} are bounded only for a Courant number up to one
(\cref{sec:disc:convection}). With either in use the solver computes, after
every step, the Courant number of the step, $\Delta t$ times the larger of
the rates of the state it starts from and the state it ends at, a rate being the
largest over the cells of $\sum_\text{axes} \max(\text{speed on the two
faces})/\Delta$:
\begin{itemize}
  \item \key{Concentration}, well-mixed: the speed is $|u|$ of the velocity
    that advects $c$, over $\Delta x$ and $\Delta y$.
  \item \key{Concentration}, \val{"CZ"}: the volume fluxes of the layers,
    $|Q|/(h\,\Delta x)$ in $x$ and $y$ and $|W|/(h\,\Delta\zeta)$ through the
    $\zeta$ faces, taken from the last assembly of the step.
  \item \key{Height}: the characteristic speed $\partial Q/\partial h$ of
    the upwinded fluxes, $3|u|$ of the in-plane gravity and the centrifugal
    term ($Q \propto h^3$) plus $2|u|$ of the Marangoni flux ($Q \propto
    h^2$), each $u$ the velocity of that term. The capillary and disjoining
    fluxes are not upwinded and are not counted.
\end{itemize}
The first step above one is reported in \file{Warnings.txt}
(\file{numerics/\bk CourantCheck.h}) --- grep for
\code{Courant number} --- and at the end of the run the number of such
steps and the largest value:
\begin{plain}
Courant number 1.43 at step 1 (t = 0.00143 s) with Numerics.Convection.Concentration = "Koren": a limited scheme keeps the concentration free of over- and undershoots only up to 1. Reported once; the count of such steps is given at the end of the run.
Courant number above 1 in 7 of 7 steps (largest 1.43) with Numerics.Convection.Concentration = "Koren".
\end{plain}
The largest value of the run is written to the log and to \file{JobInfo.txt}
(\code{Largest Courant number <conc>}, \code{<height>}) whether it passed one
or not. Nothing else follows from it: the step is neither shortened nor
rejected, and the results are those of a run without the check. To stay
below one, shorten the step: \key{Numerics.InitialDt}, or the cap
\key{Numerics.TimeStepControl.MaxDtFactor} sets on it
(\cref{sec:ref-numerics:tsc}).

\paragraph{Overshoots past $c = 1$.} The schemes that are not bounded ---
and the limited ones above a Courant number of one --- overshoot at a sharp
front of the concentration, and next to a film that has dried to $c$ near 1
the overshoot passes 1. The well-mixed model moves the solute above 1 to the
cells around (\cref{sec:disc:convection}) and keeps it; the amount moved over
the run, summed over the steps as a fraction of the initial solute, is
written to \file{JobInfo.txt} (\code{Solute moved from c > 1 [-]}) and, when
it is not zero, to the log. A front from $c = 0.99$ to 0.3 crossing a domain of
100 cells once moves 1.7 to \SI{2.3}{\percent} with UTOPIA, KK and QUICK and
\SI{4}{\percent} with Central4. It measures an error of the scheme, not of the
balance; \val{"Koren"} at a Courant number up to one moves nothing, and a
finer grid moves less. No shipped case moves any. Only a film holding more
solute than liquid in all could not take it back; that remainder is cut, reported once in \file{Warnings.txt} as
it happens --- grep for \code{cut at c = 1} --- with the total at the end of
the run, and written to \file{JobInfo.txt} (\code{Solute cut at c = 1 [-]}).

\paragraph{Undershoots below $c = 0$.} The same schemes undershoot on the
other side of a front, and next to a film with almost no solute the
undershoot passes 0. Both models fill a cell below 0 from the cells around
that hold solute (\cref{sec:disc:convection}), so $c$ stays at or above 0 and
the solute is kept; the amount moved, as a fraction of the initial solute, is
written to \file{JobInfo.txt} (\code{Solute moved to c < 0 [-]}) and, when it
is not zero, to the log. A front from $c = 0.7$ to 0.005 crossing a domain of
100 cells once moves 3.2 to \SI{4.4}{\percent} with UTOPIA, KK and QUICK and
\SI{7.7}{\percent} with Central4; \val{"Koren"} and \val{"VanLeer"} at a Courant
number up to one move nothing. No shipped case moves any. Only a film with
less than no solute in all could not be filled; the rest is made up, reported
once in \file{Warnings.txt} --- grep for \code{made up at c = 0} --- with the
total at the end of the run, and written to \file{JobInfo.txt}
(\code{Solute made up at c = 0 [-]}).

\section{The two \key{LinearSolver} blocks}
\label{sec:ref-numerics:linearsolver}

\key{Numerics.HeightEquation.LinearSolver} and
\key{Numerics.ConcEquation.LinearSolver} are both required, and both take
exactly the same five keys with the same defaults and the same allowed values;
the parser code is duplicated verbatim. You can therefore configure the two
solves independently --- a different backend for the height equation and for
the concentration equation is legitimate. \Cref{ch:linearsolver} compares the
backends and their cost.

\begin{keytable}{Keys of \key{Numerics.HeightEquation.LinearSolver} and of \key{Numerics.ConcEquation.LinearSolver}\label{tab:linsolver-keys}}\key{MaxIterations} & int & \val{5000} & Iteration cap of the BiCGStab solve, at least 1. \\
\key{Epsilon} & double & \val{1e-13} & Convergence tolerance on the (possibly scaled) relative residual, positive. For the film height of the well-mixed model it is the tightest tolerance a Newton pass is solved to: earlier passes are solved more loosely, as the iteration needs (\cref{sec:linsolve:ew}). \\
\key{RowScaling} & bool & \val{true} & Divide each row of the matrix and the right-hand side by its diagonal before solving. Does not change the solution; strongly affects conditioning. Honored by both backends. \\
\key{PreconditionType} & string & \val{"ILU"} & Preconditioner: \val{"None"}, \val{"Jacobi"} or \val{"ILU"}. Each backend maps these onto its own machinery (\cref{ch:linearsolver}); with \val{"ILU"}, a cheaper preconditioner is tried first and this one is the fall-back (\cref{sec:linsolve:eigen,sec:linsolve:conc}). \\
\key{Backend} & string & \val{"auto"} & Implementation: \val{"auto"}, \val{"eigen"} or \val{"amgcl"}. \val{"auto"} selects amgcl for a 2-D film if the binary was built with it, otherwise eigen. \val{"eigen"} is bit-reproducible regardless of thread count; \val{"amgcl"} is not. Any other value is rejected. \\
\end{keytable}


\begin{warn}[Requesting \val{"amgcl"} in a build without it falls back silently]
  A binary built without AMGCL support answers \key{Backend = "amgcl"} with a
  warning and then uses \val{"eigen"}. The run continues and the results are
  valid, but they are not the results you asked for, and comparing them
  against a run of the same input on an AMGCL build will show differences at
  the level of the linear solver's tolerance. Nothing else in the run says so
  in words; the warning in \file{Warnings.txt} is the only message about the
  substitution, though the \code{Solver Info} section of \file{JobInfo.txt}
  does record \val{"eigen"} as the backend that ran. The warning reads
  \code{LinearSolver: Backend="amgcl" was requested, but this build has no
  AMGCL. Falling back to "eigen".} Grep for \code{has no AMGCL}
  (\cref{sec:output:warnings}).
\end{warn}

\begin{note}[\file{JobInfo.txt} records which backend actually ran]
  \code{C\_ImplicitSolver::\bk BackendName()} reports the \emph{resolved}
  backend of an equation, and the model drivers write it into the
  \code{\#\#\#\#\#\# Solver Info \#\#\#\#\#\#} section of \file{JobInfo.txt}
  alongside the model name, one line for the height equation and, where there
  is one, a second for the concentration equation. Because the name is
  resolved, \val{"auto"} appears as whatever it became and a fall-back appears
  as the implementation that ran, so a result directory tells you which
  linear solver produced it. The same two lines are echoed to standard output.
  \Cref{sec:output:jobinfo} shows the section. Setting \key{Backend}
  explicitly in both input files is still the tidier way to set up a
  comparison, but it is not the only record.
\end{note}

\section{\key{Numerics.ManufacturedSolution}}
\label{sec:ref-numerics:mms}

For verification only, and optional. With \key{Enable} set, the run starts
from
\[
  h = \key{Mean} + \key{Amplitude}\,\cos(k_x x - \Omega t)\cos(k_y y),\qquad
  k = 2\pi/\key{Wavelength},
\]
over the substrate $s = \key{SubstrateMean} + \key{SubstrateAmplitude}\,
\cos(2\pi x/L_{s,x})\cos(2\pi y/L_{s,y})$ and with
the concentration $c = \key{ConcMean} + \key{ConcAmplitude}\,\cos(k_{c,x} x
- \Omega_c t + \phi_c)\cos(k_{c,y} y)$ (every $k_y = 0$ in 1-D), and the
equations get the source terms that make these their exact solution, so that
the difference between them and the computed fields is the discretization
error and nothing else. \file{tools/mms.py} uses it
(\cref{sec:verification:mms}). Every term of the well-mixed model enters the
sources; in \val{"CZ"} the concentration is the same at every height in the
film, and the evaporation and the disjoining pressure are refused; in both
models the nozzle is refused. The substrate comes from here, so a shape in \key{Substrate} is
refused as well. The film, the substrate and the concentration are set from
these formulas at the start, whatever \key{Initial} says; the edges have to be
periodic, with the domain a whole number of wavelengths across.

\begin{keytable}{Keys of \key{Numerics.ManufacturedSolution}\label{tab:mms-keys}}\key{Enable} & bool & \val{false} & Add the manufactured sources. \\
\key{Mean} & double & \val{0} & Mean film thickness in meters; greater than \key{Amplitude}. \\
\key{Amplitude} & double & \val{0} & Amplitude in meters, at least 0. \\
\key{Omega} & double & \val{0} & Angular frequency in rad/s; the wave travels in $+x$ at $\Omega/k_x$. \\
\key{Wavelength} & array & \val{[1, 1]} & Wavelengths $(L_x, L_y)$ in meters. Length exactly~2. \\
\key{ConcMean} & double & \val{0.5} & Mean concentration. \val{"CZ"} uses it only when \key{ConcAmplitude} or \key{ConcStratAmplitude} is non-zero, and otherwise keeps \key{Initial.Concentration.Value}. \\
\key{ConcAmplitude} & double & \val{0} & Its amplitude; the concentration must stay inside $(0, 1)$. In \val{"CZ"}, the same at every height in the film; \key{ConcStratAmplitude} adds a part that is not. \\
\key{ConcOmega} & double & \val{0} & Its angular frequency in rad/s. \\
\key{ConcPhase} & double & \val{0} & Its phase $\phi_c$ in radians. \\
\key{ConcWavelength} & array & \val{[1, 1]} & Its wavelengths in meters. \\
\key{ConcStratAmplitude} & double & \val{0} & \val{"CZ"} only: the amplitude $A_s$ of a part $A_s\cos(2\pi x/L_x - \Omega_s t)\cos(2\pi y/L_y)\cos\pi\zeta$ that varies through the film; $|\key{ConcAmplitude}| + |A_s|$ must keep $c$ inside $(0, 1)$. The well-mixed model refuses it. \\
\key{ConcStratOmega} & double & \val{0} & Its angular frequency $\Omega_s$. \\
\key{ConcStratWavelength} & array & \val{[1, 1]} & Its wavelengths $L_x$, $L_y$ in meters. \\
\key{SubstrateMean} & double & \val{0} & Mean substrate height in meters. \\
\key{SubstrateAmplitude} & double & \val{0} & Amplitude of the substrate relief in meters. \\
\key{SubstrateWavelength} & array & \val{[1, 1]} & Its wavelengths $(L_{s,x}, L_{s,y})$ in meters. \\
\end{keytable}

\section{\key{Restart}}
\label{sec:ref-numerics:restart}

Optional, at most once. A restart continues from a previous result instead of
building the film from \key{Initial}. Because the ASCII field file is
hand-editable, this block is also the route by which an analytically
constructed profile --- for instance one written by
\file{tools/make\_equilibrium\_1d.py} --- is injected into a run.

\begin{keytable}{Keys of \key{Restart}\label{tab:restart-keys}}\key{Enable} & bool & \val{false} & Continue from a previous result instead of the \key{Initial} block. When false, every other key here is read and ignored. \\
\key{ASCII} & bool & \val{false} & Read the ASCII neutral \file{.dat} form instead of the binary \file{.vts} form. \\
\key{FileName} & string & \val{"restartfile"} & Path of the file to read, relative to the working directory. No extension is appended, so give the full name. A file that is not there stops the run before its result directory is made. \\
\key{RestartTime} & double & \val{0} & Physical time the continued run starts at, in seconds, non-negative. Divided by $T_0$ internally. It does \textbf{not} shift \key{EndTime}. \\
\key{FileIndex} & int & \val{0} & Number the output-file counter resumes from, non-negative, so that new field files do not overwrite the old ones. \\
\end{keytable}

Three properties of the restart path are worth knowing before you rely on it.

\begin{warn}[\key{Output.Dimensioned} decides how the restart file is read]
  Nothing in \key{Restart} says what units the file is in. The thickness
  conversion factor is taken from \key{Output.Dimensioned}: with dimensioned
  output the file is assumed to hold micrometers, and with non-dimensional
  output the factor is 1. \textbf{A file written under one setting cannot be
  read back under the other}, and the mismatch is a silent factor of $h_0/10^{-6}$
  in thickness, not an error.
\end{warn}

Both readers work for both concentration models, in 1-D and 2-D, and both
verify the grid size --- the row count of a \file{field*.dat}, the point counts
in the header of a \file{.vts} --- and abort the run on a mismatch, so a
restart onto a different mesh is not possible. The reason for any failure is
written to \file{FatalErrors.txt} and to the terminal: a file of the wrong
kind (a \file{field*.dat} with \key{ASCII = false}, say) is reported as such
at once. For the well-mixed models the boundary conditions are re-applied after
reading, so whatever ghost values the file contained are overwritten; the
\val{"CZ"} \file{.vts} carries its ghost layers and they are taken as written,
and after an ASCII restart they are set by the boundary conditions at the first
solve.

A restart reproduces the state it was given, to the last digit or one unit in
it: \code{ctest} runs \code{restart\_roundtrip}, which restarts each model from
each format and compares the first field the restarted run writes with the
field it was started from, together with its time and file number; for
\val{"CZ"} it also restarts from a \file{.vts} rewritten without the
thickness and the substrate, from which the film is recovered from the point
heights.

Every field output also writes \file{field*.state} beside the field files,
with the same number. It holds what the field file does not: the time to the
bit, the step count, the time step and the one before it, the counters of the
adaptive step, the output schedule, the running totals of
\file{History.dat}, the steady-state count, the film and (in \val{"CZ"}) the
concentration of the step before, and $h$, $s$ and $c$ with their ghost
layers. When a restart finds this file beside \key{FileName}, written by the
same model at \key{RestartTime} (the time stamp in the \file{.pvd} is
enough, and so are six digits of it), the run continues exactly as if it had not been interrupted: the
same field files, \file{History.dat}, \file{CalcInfo.dat} and probe rows, and
the same time steps (\code{restart\_chain}, \cref{sec:verification:bcio}).
The state at the restart time is not written again, so \key{FileIndex}
should be the file number plus one. \file{History.dat}, \file{CalcInfo.dat},
the probes and the \file{.pvd} continue in the existing files when the
restart writes into the same directory (\key{Overwrite = true}); their rows
after the restart time are dropped, and so is a last row without its
newline, left by a run killed while writing it. The field file stays
authoritative: if it was edited, it is used as it stands and a warning says
so. A damaged or inconsistent state file is ignored with a warning, on the
terminal and in \file{Warnings.txt} (\code{... is damaged and is not used}),
and the run starts from the field file alone. A damaged field file --- a
short row in \file{field*.dat}, a block of the wrong length, a \file{.vts}
that does not end with \code{</VTKFile>} --- is refused, with the line or the
block named, rather than read partly uninitialized. The solver itself never
leaves such a file behind: every field and state file is written under a
temporary name and renamed when whole (\cref{sec:output:failures}), and a run
stopped by Ctrl-C or \code{SIGTERM} writes the state it reached, which
continues exactly.

From a field file without its state file --- for example one written by
hand --- only the field is taken. The output schedule continues on the
multiples of the interval, but the time step returns to \key{InitialDt}, the
totals of \file{History.dat} start from zero and its conservation errors are
taken against the restart state. The result then differs from an
uninterrupted run at the level of the time-stepping error with an adaptive
step, and to the digits of the file (about $10^{-13}$ for the well-mixed
model, $10^{-9}$ for \val{"CZ"}) with a fixed one.

Finally, remember that \key{EndTime} is absolute: a restart at
\key{RestartTime = 5} with \key{EndTime = 4} stops immediately.

\section{Keys the parser accepts but no case uses}
\label{sec:ref-numerics:unused}

The cross-check behind this manual, \file{tools/doc\_keys.py} (run by
\code{ctest} as \code{doc\_keys}), compares every block path and key used in
the 38 input files of \file{cases/}, \file{cases/eq/} and
\file{cases/scaling/} included, against what the parser accepts. \Cref{tab:unused-keys} lists what the parser supports
that no shipped case exercises. These keys are not deprecated and not
untested in the sense of being known-broken; they are simply outside the
coverage of \file{cases/}, so a run that uses them is exploring territory the
regression suite does not visit. Check the result rather than assuming it.

\begin{yzntable}{Parser features not exercised by any case in \file{cases/}\label{tab:unused-keys}}{@{}L{0.30\textwidth} L{0.58\textwidth}@{}}
\textsf{\bfseries Block} & \textsf{\bfseries Keys never used} \\ \midrule
root & \key{LengthScaleType}, \key{LengthScale}; every case therefore runs with the capillary length scale. \\
\key{Physics.Evaporation} & \key{PowerN}, \key{ConcR}, \key{EndConcType}, \key{EndConcValue}; the \val{"PowerN"} concentration law and the dry-out threshold are never covered. \\
\key{Physics.Slip} & The whole block; every case runs without slip. \code{analytic} (\code{leveling\_slip\_*}, \code{eq\_slip}, \code{wall\_balance\_slip}), \code{mms} (\code{wm\_slip\_1d}, \code{wm\_all\_slip\_1d}) and \code{yuzen\_tests} exercise it. \\
\key{Physics.Source} & \key{ConcType}; every nozzle case supplies its liquid at the local concentration. \code{analytic} exercises \val{"supplied"} (\code{nozzle\_supplied}, \code{nozzle\_mixing}). \\
\key{Numerics} & \key{SteadyStateCriterion}, \key{SteadyStateSteps}; the steady-state test always runs with its defaults. \\
\key{Numerics.TimeStepControl} & \key{ScaleFloor}. \\
\key{Numerics.HeightEquation} & \key{NewtonEpsilonRel}, \key{NewtonMinIterations}, \key{JacobianUpdateInterval}. \\
\key{Numerics.Convection} & \key{Concentration}; every case runs its default, \val{"UTOPIA"}. Only \file{wet\_coupled} sets \key{Height} (to \val{"UTOPIA"}); \code{ctest} exercises the other schemes (\code{advection}). \\
Both \key{LinearSolver} blocks & \key{Backend}; no case pins a backend, so all of them run whatever \val{"auto"} picks. \\
\key{Output.FieldData} & \key{Times}, \key{ExactTimes}, \key{WritePotential} (\code{analytic} and \code{restart\_roundtrip} exercise it); the listed output times are exercised through a \key{LogSchedule} (\file{leveling\_1d\_logout}), and \key{Times} by \code{input\_guards}. \\
\key{Output.Energy} & The whole block; \code{analytic} (\code{energy\_*}) and \code{restart\_roundtrip} exercise it. \\
\key{Output.ProbeList} & The entire \key{Probe} sub-block --- \key{X}, \key{Y}, \key{Z} --- although every case declares an empty \key{ProbeList} with a \key{TimeInterval}. \\
\key{Initial.SurfaceShape} & \key{RectBump}, \key{FlatPlane}, \key{MexicanHat}, \key{SmoothConvex}, \key{EdgeBump}; only \key{EllipseBump} and (in \file{leveling\_1d\_logout}) \key{PeriodicWave} are used. \\
\key{Substrate} & \key{RectBump}, \key{PeriodicWave}, \key{FlatPlane}, \key{MexicanHat}, \key{SmoothConvex}; again only \key{EllipseBump}. \\
\key{Numerics.ManufacturedSolution} & The whole block, which is for verification only; \code{mms} and \code{input\_guards} write their own inputs with it. \\
\end{yzntable}

\key{Substrate.EllipseBump.Order} is never given either --- only the
\key{Initial} copy sets it --- so the substrate-side ordering path is
untested. Two further observations come out of the same cross-check:
\key{Physics.DisjoiningPressure} appears only in \file{cases/eq/} and
\file{cases/wet\_*}, and every case supplies \key{Physics.Evaporation},
\key{Physics.ConcField}, \key{Output.ProbeList}, \key{Substrate} and
\key{Restart}, so no shipped case relies on what an absent optional block
means (\cref{ch:ref-physics}); \code{input\_guards} runs one that does.

\section{\key{InputFormatVersion}}
\label{sec:ref-numerics:formatversion}

Every shipped case file sets \key{InputFormatVersion = "1.0"} at root level.
It records which version of the input format the file was written for. It is
optional, since there has only been one, and when it is present it must read
\val{"1.0"}: a file written for a later format is refused rather than
half-understood.

The other version gate is \key{SolverName}, which must read exactly
\val{"Yuzen-1.0"}. It names the version \emph{series}, \code{MAJOR.MINOR}, and
not the full version of the binary: the solver builds the accepted string from
\code{C\_Yuzen::\bk VersionSeries()}, so a patch release goes on reading the
same input files unchanged. \Cref{ch:policy} sets out what that promises.
